% chapters/99_appendix.tex

\chapter*{Apêndice: O Manual do Arquiteto}
\addcontentsline{toc}{chapter}{Apêndice: O Manual do Arquiteto}

\section*{Uma Nota do Autor}
Esta história nasceu nas ruas cinzentas e escorregadias de São Paulo, mas suas raízes são globais. Como Elara, passei anos olhando para a ``Garoa Residual'' --- os sistemas opacos, ineficientes e muitas vezes cruéis do Mercado Legado --- e me perguntando se poderíamos simplesmente... projetar algo melhor.

\textit{System Call} é mais do que um romance. É uma especificação funcional para uma transição. A São Paulo sobre a qual você lê é um espelho da nossa: um mundo onde somos cada vez mais tratados como variáveis em um loop de maximização de lucros. Mas o Coletivo também é real. Ele existe onde quer que as pessoas escolham a utilidade em vez da arbitragem, e a comunidade em vez da dívida.

Eu moro perto das serras onde a Garoa Residual é mais espessa, e estou escrevendo o código para a primeira Malha. Este livro é o meu convite para você se juntar à equipe de design.

\section*{Glossário do Coletivo}
\begin{description}[style=nextline, leftmargin=1em, font=\bfseries\headerfont\color{collectiveBlue}]
    \item[Corpo de Conhecimento (BoK)] O registro global, verificado por pares, de habilidades e definições de tarefas. O ``código-fonte'' da capacidade humana.
    \item[Fundo Comum] A interface entre a Malha e o Mercado Legado, usada para arrendar recursos externos até que possam ser internalizados.
    \item[Dissociação] O ato de remover um ativo físico ou digital do registro legal e financeiro do Legado e ancorá-lo em um Nó Sociedade.
    \item[Sangria] O redirecionamento estratégico de recursos para longe dos ciclos de capital e para dentro de loops metabólicos autossustentáveis.
    \item[Nós de Serviço] Unidades modulares de produção (ex.: Comida, Remédios, Energia) definidas por uma WBS específica.
    \item[Nós Sociedade] Nós de Serviço específicos para governança do aspecto social e não produtivo da Malha, as entidades ``legais'' do Coletivo.
    \item[Valor] Uma unidade de conta não monetária que representa a contribuição de um membro para a utilidade da malha.
\end{description}

\section*{O Repositório da Realidade}
A lógica descrita neste livro --- as Estruturas de Divisão de Trabalho (WBS), as especificações técnicas para os Nós de Serviço, e o White Paper detalhando a Guerra Metabólica --- não está oculta. Ela é de código aberto e está pronta para auditoria.

Se você é um Designer, um Desenvolvedor, ou simplesmente uma Variável pronta para se tornar um Arquiteto, você pode acessar a lógica central do Coletivo aqui:

\begin{center}
    \headerfont\color{collectiveBlue} \url{https://github.com/raffaine/collective}
\end{center}

Este repositório contém as plantas do sistema que Elara descobriu. Ele inclui as especificações técnicas para a Malha, a estrutura do BoK e os rascunhos dos protocolos para a Ancoragem de Sociedade. É o começo da Dissociação.

\section*{Sua Primeira System Call}
Se você deseja iniciar seu processo de integração, não procure por um site ou um aplicativo. Procure pelo \textbf{Sinal}.
\begin{enumerate}
    \item \textbf{Identifique a Demanda Latente:} Encontre uma coisa na sua vizinhança imediata que esteja quebrada, sem uso ou que seja necessária, mas que o mercado recusa consertar por não ser ``lucrativa''.
    \item \textbf{Modele a Utilidade:} Quantas pessoas são afetadas? Quais habilidades são necessárias para consertar isso? Este é o começo da sua primeira \textbf{WBS}.
    \item \textbf{Forme um Loop de Pares:} Encontre uma outra pessoa que compartilhe da mesma necessidade. Vocês agora são uma \textbf{Micro-Célula}.
\end{enumerate}

\begin{center}
    \textit{A Garoa Residual só existe enquanto concordarmos em caminhar por ela às cegas.} \\
    \textbf{Abra os olhos. Responda ao chamado.}
\end{center}
